This section discusses safety, security, and safeguards, and how each provides motivation for this work.
Safety and security are considered in the context of potential transient events and their consequences.
Security is examined from the perspective of potential insider sabotage, which partially overlaps with safety analysis. 
A saboteur could initiate an event or combination of events that threaten the general public, serving as a useful distraction for hostile agents.
Finally, safeguards analyses are useful for determining if a diversion of \gls{SNM} has occurred, and if so, how much and the category of the diverted material.
One diversion fingerprinting technique uses a kinetics response, since diversion of fissile material alters the delayed neutron response to a transient.
By improving the \gls{DNP} group parameters for \glspl{MSR}, previously undetectable or uncertain diversions will become more readily identifiable.
Section \ref{sec:transimpact} provides detail on what analyses will be conducted for safety, security, and safeguards.